% !TEX program = xelatex
%% ============================================================
%% 《中文信息学报》投稿 LaTeX 模板（指导增强版）
%% Journal of Chinese Information Processing (JCIP)
%%
%% 说明：
%% 1. 本模板在可编译版本基础上，吸收了你提供的 Word/PDF 指导模板中的关键写作说明。
%% 2. 这些说明主要以"源码注释"的形式放在对应位置，便于边写边看；默认不会显示在 PDF 中。
%% 3. 如需最终投稿，请删除示例文字并替换为真实内容。
%% 4. Overleaf 编译器请设置为 XeLaTeX。
%% ============================================================

\documentclass[10pt,twoside,twocolumn]{article}

% ── 中文与字体 ─────────────────────────────────────────────
\usepackage{ctex}
\usepackage{fontspec}
\IfFontExistsTF{Noto Serif CJK SC}{%
  \setCJKmainfont{Noto Serif CJK SC}
  \setCJKsansfont{Noto Sans CJK SC}
  \setCJKmonofont{Noto Sans CJK SC}
}{}

% ── 页面与基础宏包 ───────────────────────────────────────
\usepackage[
  a4paper,
  top=28mm,
  bottom=25mm,
  left=19mm,
  right=19mm,
  columnsep=6mm,
  headsep=6mm,
  headheight=24pt,
  footskip=10mm
]{geometry}

\usepackage{amsmath,amssymb,bm}
\usepackage{graphicx}
\usepackage{booktabs}
\usepackage{caption}
\usepackage{fancyhdr}
\usepackage{titlesec}
\usepackage{indentfirst}
\usepackage{microtype}
\usepackage{cite}
\usepackage[hidelinks]{hyperref}
\usepackage{xcolor}

\captionsetup{font=small,labelfont=bf,labelsep=space}
\captionsetup[figure]{position=below,name=图}
\captionsetup[table]{position=above,name=表}
\setlength{\parindent}{2em}

% ── 期刊信息与眉题 ───────────────────────────────────────
\newcommand{\jcipvol}{**}
\newcommand{\jcipno}{*}
\newcommand{\jcipyear}{202*}
\newcommand{\jcipmonth}{*}
\newcommand{\articleno}{1003-0077（2017）00-0000-00}

\newcommand{\runningtitletext}{请在此填写文章眉题（短标题）}
\newcommand{\runningtitle}[1]{\gdef\runningtitletext{#1}}

% ── 页眉页脚 ─────────────────────────────────────────────
\pagestyle{fancy}
\fancyhf{}
\fancyhead[LE]{\zihao{-5}\thepage\quad 中文信息学报\quad 第\jcipvol 卷}
\fancyhead[CE]{\zihao{-5}\runningtitletext}
\fancyhead[RO]{\zihao{-5}第\jcipvol 卷\ 第\jcipno 期\quad 中文信息学报\quad Vol.\jcipvol，No.\jcipno}
\fancyhead[CO]{\zihao{-5}\runningtitletext}
\fancyfoot[LE,RO]{\zihao{-5}\thepage}
\renewcommand{\headrulewidth}{0.4pt}

\fancypagestyle{jcipfirst}{%
  \fancyhf{}
  \fancyhead[L]{\zihao{-5}
    \begin{tabular}[b]{@{}l@{}}
      第\jcipvol 卷\ 第\jcipno 期\\
      \jcipyear 年\ \jcipmonth 月
    \end{tabular}}
  \fancyhead[C]{\zihao{-5}\textbf{中文信息学报}\\[-1pt]
    \small JOURNAL OF CHINESE INFORMATION PROCESSING}
  \fancyhead[R]{\zihao{-5}
    \begin{tabular}[b]{@{}r@{}}
      Vol.\jcipvol，No.\jcipno\\
      \jcipmonth，\jcipyear
    \end{tabular}}
  \fancyfoot[C]{\zihao{6}文章编号：\articleno}
  \renewcommand{\headrulewidth}{0.4pt}
}

% ── 标题样式 ─────────────────────────────────────────────
% 一级标题：小四黑体；"引言、前言、概论"等以"0"为章节编号
% 二级标题：五号黑体
% 三级标题：五号宋体
\titleformat{\section}{\zihao{-4}\heiti}{\thesection}{1em}{}
\titleformat{\subsection}{\zihao{5}\heiti}{\thesubsection}{1em}{}
\titleformat{\subsubsection}{\zihao{5}\songti}{\thesubsubsection}{1em}{}
\setcounter{secnumdepth}{3}

% ── 摘要、关键词、中图分类号 ─────────────────────────────
\renewenvironment{abstract}{%
  \noindent{\zihao{-5}\heiti 摘要：}\zihao{-5}\kaishu\ignorespaces
}{\par\medskip}

\newenvironment{enabstract}{%
  \noindent{\zihao{-5}\textbf{Abstract: }}\zihao{-5}\ignorespaces
}{\par\medskip}

\newcommand{\keywords}[1]{%
  \noindent{\zihao{-5}\heiti 关键词：}{\zihao{-5}\kaishu #1}\par\smallskip}
\newcommand{\enkeywords}[1]{%
  \noindent{\zihao{-5}\textbf{Key words: }}{\zihao{-5} #1}\par\medskip}
\newcommand{\clccode}[2]{%
  \noindent{\zihao{-5}\heiti 中图分类号：}{\zihao{-5}#1}%
  \hfill{\zihao{-5}\heiti 文献标识码：}{\zihao{-5}#2}\par\bigskip}

% ── 首页脚注信息 ─────────────────────────────────────────
\newcommand{\received}[2]{%
  \footnotetext{\zihao{6}{\heiti 收稿日期：}#1；{\heiti 定稿日期：}#2}}
\newcommand{\fundinfo}[1]{%
  \footnotetext{\zihao{6}{\heiti 基金项目：}#1}}
\newcommand{\corrmark}{\textsuperscript{\ensuremath{\dagger}}}
\newcommand{\corrauthornote}[2]{%
  \footnotetext{\zihao{6}$^{\dagger}$\,通信作者：#1，E-mail: \href{mailto:#2}{#2}}}

% ── 作者简介与照片 ───────────────────────────────────────
\newlength{\photow}
\newlength{\photoh}
\setlength{\photow}{1.8cm}
\setlength{\photoh}{2.4cm}

\newcommand{\authorentry}[2]{%
  \par\noindent
  \begin{minipage}[t]{\photow}
    \if\relax\detokenize{#1}\relax
      \fbox{\rule{0pt}{\photoh}\rule{\dimexpr\photow-2\fboxsep-2\fboxrule\relax}{0pt}}%
    \else
      \includegraphics[width=\photow,height=\photoh]{#1}%
    \fi
  \end{minipage}%
  \hspace{4pt}%
  \begin{minipage}[t]{\dimexpr\linewidth-\photow-4pt\relax}
    \zihao{6}#2
  \end{minipage}\par\smallskip
}

% ── 数学函数名正体 ───────────────────────────────────────
\newcommand{\func}[1]{\mathrm{#1}}

\begin{document}
\thispagestyle{jcipfirst}
\runningtitle{请在此填写文章眉题（短标题）}

% ─────────────────────────────────────────────────────────
% 首页题头区
% 中文标题：三号
% 中文作者：四号宋体，单位对应关系用右上角上标标注
% 单位：小五号
% 英文标题：四号
% 英文作者：五号
% 英文单位：小五号
% ─────────────────────────────────────────────────────────
\twocolumn[{%
  \begin{center}
    {\zihao{3}\heiti 中文标题题目三号}\par\medskip

    {\zihao{4}\songti
      作者 1$^{1}$\quad 作者 2$^{1,2}$\quad 作者 3$^{1,2}$
      \quad {\zihao{4}\songti 作者四号宋体，作者和单位的对应关系标注在作者姓名的右上角}
    }\par\smallskip

    {\zihao{-5}\songti
      （1. 一级单位名称\ 二级单位名称，省\ 市\ 邮编；
      2. 一级单位名称\ 二级单位名称，省\ 市\ 邮编）{\zihao{6}单位小五号}
    }\par\medskip
  \end{center}

  \begin{abstract}
中文摘要 100--300 字，用"该文"不用"本文"。概括论文内容，写明研究目的、方法、结果和结论，
文摘要具体化，如结果的百分比等。文摘要开门见山，只叙述新信息和发现，而不必写课题研究的背景信息和
的研究细节。文摘要表明作者原创性工作，突出要点。文摘中只要最关键的数据。文摘中不能出现图、表、参
考文献等数据。文摘中不表述个人观点，不出现未来计划。文摘中的缩写要有全称，专业词汇准确。不说无用
的话，不需要自己标榜自己的研究结果，避免类似下列的句子该出现在文摘中："本文的有关研究工作是对以
往工作的一个极大的改进"，"本工作首次实现了……"，"经检索尚未发现与本文类似的文献"等。具有自我独
立性。文摘第一句应避免与题目（Title）重复。文摘中应避免出现特殊字符，即各种数学符号、上下脚标及
希腊字母，相应内容改用文字表达或文字叙述。
  \end{abstract}
  \keywords{三个左右为宜，用分号；分开}
  \clccode{TP391}{A}

  \begin{center}
    {\zihao{4}\textit{English title Title 四号}}\par\medskip

    {\zihao{5}
      Author1$^{1}$,\ author1$^{1,2}$,\ and author1$^{1,2}$
      \quad {\zihao{5} Name 五号}
    }\par\smallskip

    {\zihao{-5}
      (1. Unit, city, province zip code, China;\ 2. Unit, city, province zip code, China)
      {\zihao{-5} Depart.Correspond 小五号}
    }\par\medskip
  \end{center}

  \begin{enabstract}
与中文文摘相对应，英文文摘应符合英文语法，概括论文内容，写明研究目的、方法、结果和结论。尽量
用短句子并避免句型单调。用过去时态叙述作者工作，用现在时态叙述作者结论。避免以下字句出现：如
``It is reported \ldots{}'' ``Extensive investigations show that\ldots{}'' ``The author discusses \ldots{}''
``This paper concerned with \ldots{}''；文摘开头的 ``In this paper,''。一些不必要的修饰词，如
``in detail''、``briefly''、``here''、``new''、``mainly'' 也尽量不要。能用名词做定语不要用动名词做
定语，能用形容词做定语就不要用名词做定语。例如：用 selection method 不用 selecting method 用
experimental results 不用 experiment results，可直接用名词或名词短语作定语的情况下，要少用 of
句型。例如 用 selection method 不用 method of selection 用 camera curtain shutter 不用 curtain
shutter of camera 用 sentence structure 不用 structure of sentence 注意冠词用法，不要误用，滥用或
随便省略冠词。避免使用一长串形容词或名词来修饰名词，可以将这些词分成几个前置短语，用连字符连接名
词组，作为一个形容词。尽量用主动语态代替被动语态。尽量用简短、词义清楚并为人熟知的词。慎用行话和
俗语。文词要纯朴无华，不用多姿多态的文学性描述手法。删繁从简。如用 increased 代替 has been found
to increase 用 the results show 代替 from the experimental results, it can be concluded that
文摘词语拼写，用英美拼法都可以，但在每篇文章中须保持一致。注意中英文不同的表达方法。不要简单地逐
字直译。例如，不要将 because 放在句首表达"因为"这一概念，because 表示原因语气用在文摘中过强。不
要用 ``XX are analyzed and studied(discussed)''。来直译"分析研究（讨论）"这一中文概念。用 ``XX are
analyzed'' 就可以了。尽量不要使用 not only\ldots{}but also 直译中文"不但"\ldots{}"而且"这一概念，用
and 就行了。
  \end{enabstract}
  \enkeywords{三个左右为宜，用分号；分开 Key words 五号}

  \bigskip
}]

\received{20XX-XX-XX}{20XX-XX-XX}
\fundinfo{基金名（基金号）；基金名（基金号）
{\zihao{6} 六号，核实准确完整的基金名称，如自然科学基金（基金号）}}
\corrauthornote{作者丙}{corrauthor@university.edu.cn}

\zihao{5}\songti
\setcounter{section}{-1}

%% 正文五号宋体，双栏排版
\section{标题\quad {\zihao{-4}\heiti 一级标题小四黑体，引言、前言、概论等以"0"为章节编号}}

论文需与中文信息学处理主题切合，条理清晰，主次分明，且在理论、方法或资源上有明确的贡献。

引言中，需对论文的动机、主要创新点做出清晰易懂的说明，对论文的大致结构进行简要的介绍。

相关研究中，需要对相关研究、方法做充足的调研，亦采用归类总结的方式，而不是简单的罗列。简要说明各类方法的优缺点。对公共所知的内容应尽量避免长篇大论。与本文密切相关的研究方法务必做出简要对比说明。

实验部分需采用公开的数据集，如在自建数据集上进行实验，则需要公开数据集（训练集需要全部公开，测试集可选择性公开一部分）。

对于所采用的方法的关键环节需要进行详细、清晰地描述。对于众所周知的部分，尽量不要描述地过多。模型图应自己绘制，不可直接从他人论文中截取。

对比实验的选取上，需选择近三年的方法进行对比，过于陈旧的或相关性不强的方法可不列举。同时进行必要的消融实验。

实验分析中，需清晰地说明本文模型与其他模型相比的优势及其原因，以及各模块的具体贡献及原因。

总结及展望部分需对本文方法进行简要总结，突出论文的创新点，同时对未来进一步的改进或研究方向提出明确可行的期许。

正文用"本文"，务必避免过多第一人称的表述。篇幅在 8000 字左右为宜，参考文献按文中出现顺序引用。\cite{ref1} 参考文献顺序标引

正文中，需特别注意术语和专有名词的使用。务必避免中英文混杂的情况出现，文中首次出现的术语、专有名词等，需要优先给出中文名称，如为外文翻译过来的，则需要在中文翻译名称后的括号中给出全称和缩写，如长短时记忆网络(Long Short Term Memory Network, LSTM)。

正文中，图表须注明中文图题和表题，且在正文中应明确提及。其中图的编号和图题置于图下方的居中位置，表的编号和表题置于表上方的居中位置。

示图使用黑白绘图，对于不同的内容请采用清晰的标识，以免黑白印刷造成的混淆不清。请确保图表中文字清晰。

公式用 5 号字体，其中，变量的符号应采用斜体，向量、矢量、张量、矩阵用黑斜体表示。函数（单词）用正体小写，第一个字母小写；单个字母斜体，公式采用的符号、字母表示的涵义需要说明，且需要避免一篇文章中重复使用同一个字母或符号表示不同的涵义，以免引起歧义。

\section{标题}

\subsection{标题\quad {\zihao{5}\heiti 二级标题五号黑体}}

正文

\subsubsection{标题\quad {\zihao{5}\songti 三级标题五号宋体}}

正文

% ─────────────────────────────────────────────────────────
% 参考文献
% ─────────────────────────────────────────────────────────
\begin{thebibliography}{99}
\zihao{-5}
% 参考文献小五号，只列举最主要的，必须是公开发表的书刊才能列入，最少不得少于 5 条。
% 文献按文章中出现先后顺序排列
% （各类文献严格按照主页上的《参考文献规范》）

\bibitem{ref1}
专著：作者. 题名[M]. 出版地：出版者，出版年：起止页码.

\bibitem{ref2}
期刊：作者（多作者用逗号分开，超过 3 个者用"，等"代替）. 文章题目[J]. 刊物名称，年，卷(期)：起止页码.

\bibitem{ref3}
会议论文集：作者. 题名[C]//编者. 论文集名. 出版地：出版者，出版年：起止页码.

\bibitem{ref4}
英文会议：作者. 题名[C]//Proceedings of 会议名称. 出版地：出版者，出版年：起止页码.

\bibitem{ref5}
学位论文：作者. 题名[D]. 保存单位 XX 学位论文，年份.

\bibitem{ref6}
报告：作者. 题名[R]. 保存地点：保存单位，年份.

\bibitem{ref7}
报纸文章：作者. 题名[N]. 报纸名，出版日期（版次）.

\bibitem{ref8}
标准：标准编号，标准名称[S]. 出版地：出版者，出版年.

\bibitem{ref9}
专利：专利所有者. 专利题名[P]. 专利国别：专利号，公开日期.

\bibitem{ref10}
电子文献：主要责任者. 电子文献题名[电子文献标识/载体类型].
[发表或更新日期]. 电子文献的出处或可获得地址.

\bibitem{ref11}
电子文献标识：[DB]--数据库\quad [CP]--计算机程序\quad [EB]--电子公告

\bibitem{ref12}
电子文献载体类型：[OL]--联机网络\quad [MT]--磁带\quad [DK]--磁盘\quad [CD]--光盘

\end{thebibliography}

\vspace{6pt}
\noindent{\zihao{6}\heiti 作者简介：}\par\smallskip

% 作者介绍六号，彩色正面免冠照片，本刊仅列出前三位作者信息
% 如需标出通信作者，可以为第一作者，第二作者，通信作者（在学历前加"通信作者，"）。
% 介绍包括姓名，出生年，学历，职称和主要研究领域。
% 如需标明通信作者，请保证其至少位于前三名作者之中，否则在文章中无法体现。
\authorentry{}{%
  第一作者姓名（出生年—），学历，职称，主要研究领域为。\\
  E-mail：***@***}

\authorentry{}{%
  第二作者姓名（出生年—），学历，职称，主要研究领域为。\\
  E-mail：***@***}

\authorentry{}{%
  第三作者姓名（出生年—），学历，职称，主要研究领域为。\\
  E-mail：***@***}

\end{document}
